%% =========================================================================
\section{Spectral and Variance Routes}
\label{sec:spectral-routes}
%% =========================================================================

The ensemble framework of Section~\ref{sec:ensemble} located the orbit
hypothesis CME within the reduction landscape.  This section collects the routes that
operate in the \emph{character spectrum} of the walk itself: the Weyl
hitting bridge, which converts walk character cancellation into
cofinal hitting for arbitrary squarefree starting points; the spectral
identity, which equates subquadratic visit energy~(SVE) with visit
equidistribution~(VE); the Lyapunov reformulation of the walk's
self-correcting feedback; and the tail window decomposition, which
attacks the standard orbit through temporal decorrelation of
character-sum windows.  Alongside these spectral routes we develop the
\emph{variance} routes---the per-chi cancellation bridge and the JSE
chain---which share the same energy bookkeeping.

\begin{roadmap}
\begin{enumerate}[nosep,leftmargin=*]
\item[\S\ref{sec:weyl-bridge}] Generalized MC and the Weyl hitting bridge --- character cancel $\Rightarrow$ walk hitting via strong induction; per-chi, JSE, and variance routes.
\item[\S\ref{sec:visit-eqd}] Visit equidistribution and the spectral identity --- SVE $\leftrightarrow$ VE via Parseval.
\item[\S\ref{sec:self-correcting}] Self-correcting feedback --- the Lyapunov telescope, SCD, and drift bounds.
\item[\S\ref{sec:twd}] Tail window decorrelation --- temporal windows of the standard orbit.
\end{enumerate}
\end{roadmap}

\subsection{Generalized MC and the Weyl Hitting Bridge}
\label{sec:gen-mc}
\label{sec:weyl-bridge}

The ensemble framework of Section~\ref{sec:ensemble} requires proving
MC not just for $n = 2$ but for arbitrary squarefree starting points.
This motivates two key results.

\begin{definition}[\defname{GenMullinConjecture}$(n)$]
Every prime appears in the generalized EM sequence from starting
point~$n$: $\forall\, q$ prime, $\exists\, k$,
$\seq{n}(k) = q$.
\end{definition}

\begin{theorem}[\lean{EM/Ensemble/WeylChain.lean}{gen\_hitting\_implies\_gen\_mc\_proved}]
\label{thm:gen-mc}
For squarefree $n \geq 1$: if the generalized walk hits~$-1$
cofinally for every prime~$q$,
then $\mathrm{GenMullinConjecture}(n)$ holds.
\end{theorem}

\begin{proof}
Strong induction on~$q$.  Assume every prime $p < q$ appears as
$\seq{n}(k_p)$ for some~$k_p$.  Collect all witnesses into
a uniform bound~$N$.  Cofinal hitting at~$q$ provides $k \geq N$ with
$q \mid \Prod{n}(k) + 1$.  Since all primes $p < q$ already appeared
at steps $< k$, they divide $\Prod{n}(k)$ and hence cannot divide
$\Prod{n}(k) + 1$.  So $q = \minFac(\Prod{n}(k) + 1) =
\seq{n}(k)$.

The argument is a generalized \code{captures\_target} lemma
(\lean{EM/Ensemble/WeylChain.lean}{gen\_captures\_target}), paralleling the
standard version from \S\ref{sec:bootstrap}.
\end{proof}

The connection between character cancellation and hitting uses the Weyl
equidistribution criterion, adapted to the multiplicative setting:

\begin{theorem}[\defname{WeylHittingBridge} ---
  \lean{EM/Ensemble/WeylChain.lean}{weyl\_hitting\_bridge\_proved}]
\label{thm:weyl-bridge}
If all non-trivial walk character sums satisfy
$\|\sum_{k<K} \chi(\walkZ{n}{q}{k})\|^2 \leq (\varepsilon K)^2$
for large~$K$, then either $q$ appears in $\seq{n}(\cdot)$
or the walk hits~$-1$ cofinally.
\end{theorem}

\begin{proof}
By contradiction using a \emph{Weyl test function}.  Suppose $q$ never
appears and the walk stops hitting~$-1$ past some step~$N_0$.  Then the
hit count $H(K) = |\{k < K : \Prod{n}(k) \equiv -1 \pmod{q}\}|$ is
bounded by the constant $C = H(N_0)$ for all~$K$.

Define a centered test function $\chi_{\mathrm{test}} : \NN \to \mathbb{C}$:
\[
  \chi_{\mathrm{test}}(x) \;=\;
  \begin{cases}
    0 & \text{if } x \equiv 0 \pmod{q}, \\
    (q{-}2)/(q{-}1) & \text{if } x \equiv -1 \pmod{q}, \\
    -1/(q{-}1) & \text{otherwise.}
  \end{cases}
\]
The test function satisfies $\|\chi_{\mathrm{test}}\|^2 \leq 1$,
$\chi_{\mathrm{test}}(0) = 0$, and
$\sum_{a \in \ZZ/q\ZZ} \chi_{\mathrm{test}}(a) = 0$ (the sum of
$(q{-}2)/(q{-}1)$ from $a = q{-}1$ and $(q{-}2) \cdot (-1/(q{-}1))$
from the remaining units is zero).

The partial sum of the test function along the walk telescopes to
$\sum_{k<K} \chi_{\mathrm{test}}(\Prod{n}(k)) = H(K) \cdot
(q{-}2)/(q{-}1) + (K - H(K)) \cdot (-1/(q{-}1))
= H(K) - K/(q{-}1)$.
So the walk character energy equals
\[
  E_{\mathrm{test}}(K) \;=\; \bigl(H(K) - K/(q{-}1)\bigr)^2.
\]
Since $H(K) \leq C$ while $K/(q{-}1) \to \infty$, for large~$K$ we have
$E_{\mathrm{test}}(K) \geq (K/(q{-}1) - C)^2$.  Setting
$\varepsilon = 1/(2(q{-}1))$, the hypothesis gives
$E_{\mathrm{test}}(K) \leq (K/(2(q{-}1)))^2$.  Taking square roots:
$K/(q{-}1) - C \leq K/(2(q{-}1))$, hence $K/(2(q{-}1)) \leq C$,
so $K \leq 2C(q{-}1)$.  But $K$ was chosen larger than $2Cq + 1$:
contradiction.
\end{proof}

\begin{theorem}[\lean{EM/Ensemble/WeylChain.lean}{cancel\_weyl\_implies\_mc}]
\label{thm:cancel-weyl-mc}
Walk character cancellation at $n = 2$ implies MC.
\end{theorem}

\begin{proof}
For $q = 2$: $\seq{2}(0) = 2$.  For $q \geq 3$: $q$ does not divide~$2$,
so the Weyl Hitting Bridge gives either $q \in \seq{2}(\cdot)$
(hence $q \in \{\seq{2}(k)\}$) or cofinal hitting.  Cofinal hitting gives
$q \in \seq{2}(\cdot)$ by Theorem~\ref{thm:gen-mc}, noting
that all generalized primes from $n = 2$ are odd (since $2 \mid \Prod{2}(k)$
implies $\Prod{2}(k) + 1$ is odd), so the induction only requires
odd primes $< q$, all of which have already appeared by the
inductive hypothesis.
\end{proof}

\paragraph{The Per-Chi Cancellation Bridge}
\label{sec:per-chi}

The decorrelation chain of \S\ref{sec:decorr-chain} works with
\emph{universal} StepDecorrelation (quantified over all characters).
For the JSE route, only \emph{per-chi} decorrelation is available (for
each individual non-trivial character).  The PerChiCancellationBridge
shows the same three-step chain works per-chi:

\begin{theorem}[\lean{EM/Ensemble/WeylChain.lean}{per\_chi\_cancellation\_bridge\_proved}]
\label{thm:per-chi}
\defname{PerChiCancellationBridge}: for each prime~$q$ and
character~$\chi$ with $\|\chi\|^2 \leq 1$, if the cross-terms
$\mathbb{E}^{\mathrm{sqf}}[\mathrm{Re}(\chi(m_j) \cdot
\overline{\chi(m_k)})] \to 0$ for all $j < k$, then the density of
squarefree~$n$ with $E(n, K) > (\varepsilon K)^2$ for all~$K$ tends
to zero.
\end{theorem}

\begin{proof}
The proof is identical to the universal chain (Theorem~\ref{thm:variance-bound}
$\to$ concentration $\to$ cancellation), with the StepDecorrelation
hypothesis restricted to the fixed character~$\chi$.  Each step
(variance with $C = 2$, Markov concentration, squeeze to zero) goes
through unchanged because none requires decorrelation of \emph{other}
characters.
\end{proof}

\paragraph{The Ensemble Chain in Detail}
\label{sec:ensemble-chain}

The original ensemble assembly used the following master theorem, now
archived --- two of its four hypotheses (SRE and CRTPropagationStep)
are RED-false by the absorption mechanism (Dead End~\#137), so the
theorem is a valid conditional whose antecedents fail; the \emph{live}
assembly is the orbit hypothesis CME of Section~\ref{sec:ensemble}:

\begin{theorem}[\archived{EM/Archive/Ensemble/PTArchive.lean}{ensemble\_pt\_master\_simplified} \superseded]
Under four hypotheses (SRE, CRTPropagationStep,
AccumEquidistImpliesMultEquidist, EnsembleEquidistImpliesDecorrelation),
character sums cancel for almost all squarefree starting points.
\end{theorem}

\noindent The chain:
\[
  \underbrace{\mathrm{SRE} + \mathrm{CRT} + \mathrm{Bridge}}_{\text{ANT + CRT structure}}
  \;\xrightarrow{\text{proved}}\;
  \mathrm{EME}
  \;\xrightarrow{\text{open}}\;
  \mathrm{SD}
  \;\xrightarrow{\text{proved}}\;
  \text{cancellation for a.a.~}n.
\]
The first arrow composes three proved reductions.  The second arrow
(EME~$\Rightarrow$~SD) is the sole remaining open hypothesis.  The
third arrow is Theorem~\ref{thm:sd-cancel}.

\paragraph{The JSE route.}
An alternative chain uses \emph{JointStepEquidist}~(JSE): the joint
distribution of $(\seq{n}(j), \seq{n}(k))$
is asymptotically uniform over non-zero pairs.  The JSE base case
($j = 0, k = 1$) is proved via the tower rule
(\lean{EM/Ensemble/BaseCase.lean}{jse\_base\_case}):
$P(A, B) = P(A \mid B) \cdot P(B)$, decomposing the joint density
into conditional and marginal equidistribution.

\begin{theorem}[\lean{EM/Ensemble/WeylChain.lean}{jse\_transfer\_implies\_mc}]
$\mathrm{JSE} + \mathrm{MultCancelToWalkCancel}
\;\Longrightarrow\; \MC$.
\end{theorem}

\noindent The composition:
\begin{enumerate}[nosep]
\item JSE gives per-chi decorrelation for non-trivial~$\chi$
  (\lean{EM/Ensemble/PT.lean}{joint\_step\_equidist\_implies\_step\_decorrelation}).
\item PerChiCancellationBridge (Theorem~\ref{thm:per-chi}) gives
  multiplier character cancellation for almost all squarefree~$n$.
\item \textbf{MultCancelToWalkCancel} (open): multiplier energy
  $o(K^2)$ implies walk energy $o(K^2)$.
\item WeylHittingBridge (Theorem~\ref{thm:weyl-bridge}) converts walk
  character cancellation into cofinal hitting.
\item Generalized MC (Theorem~\ref{thm:gen-mc}) and standard-EM
  bridge (Theorem~\ref{thm:cancel-weyl-mc}) give MC\@.
\end{enumerate}

The two remaining open hypotheses are:
\begin{itemize}[nosep]
\item \textbf{JointStepEquidist}: requires analytic number theory
  (Siegel--Walfisz or Bombieri--Vinogradov type estimates).
\item \textbf{MultCancelToWalkCancel}: multiplier character cancellation
  implies walk character cancellation.  Dead End~\#117 shows this is
  equivalent to CCSB/CME in difficulty: an explicit counterexample
  (alternating multipliers $\{2, 3\} \bmod 5$ with $\chi(2) = i$) gives
  multiplier cancellation yet $|W_K| = \Theta(K)$, and this
  counterexample is compatible with all EM structural properties
  (squarefree accumulator, coprime cascade, CRT invariance).
\end{itemize}

\noindent MultCancelToWalkCancel is the Marginal/Joint Barrier
(\S\ref{sec:oracle}) in ensemble form: the multiplier character sum is a
\emph{marginal} statistic (it depends on the individual multiplier
values), while the walk character sum is a \emph{joint} statistic (it
depends on all previous multipliers through the partial product).
Cancellation of individual terms does not imply cancellation of their
cumulative products.  This is why the orbit hypothesis is stated as
CME (Section~\ref{sec:dsl-route}): CME implies walk character
cancellation directly, without ever needing the multiplier-to-walk
transfer.

\paragraph{The variance route (retired).}
\label{sec:variance-route}
An earlier draft developed a second-moment route
$\mathrm{PCV} \Rightarrow \mathrm{SMB} \Rightarrow \text{concentration}
\Rightarrow \mathrm{DSL}$ (population character energy, cross-term pairs,
SD~$\Rightarrow$~PCV, CSVB~$\Rightarrow$~SMB).  Its endpoint DSL is
vacuous (Dead End~\#160), and its uncentered second-moment hypotheses
are false at the trivial character (Dead End~\#156); the framework is
archived in \code{EM/Archive/Reduction/DSLVarianceArchive.lean} (local archive).  What
survives of it is the observation that concentration gives cancellation
for \emph{almost all} squarefree starting points
(Theorem~\ref{thm:sd-cancel}) while CME asks for it at $n = 2$: a
density-one set of good starting points can exclude any particular~$n$,
the orbit-specificity barrier in its sharpest form.

\subsection{Visit Equidistribution and the Spectral Identity}
\label{sec:visit-eqd}

The ensemble framework connects to the spectral route via a
fundamental identity relating character energy to the geometry of the
walk.  For a walk of length~$N$ on $(\ZZ/q\ZZ)^\times$, define the
\emph{visit count} $V_N(a) = |\{n < N : \walkZ{2}{q}{n} = a\}|$.

\begin{theorem}[\lean{EM/Reduction/VisitEquidist.lean}{excessEnergy\_eq\_visit\_deviation}]
\label{thm:excess-energy}
The excess energy $E_\chi(N) - N$ equals $(q{-}1)$ times the
squared deviation of visit counts from the uniform distribution:
\[
  \sum_{\chi \neq \chi_0} \bigl\|S_\chi(N)\bigr\|^2
  \;=\; (q{-}1) \sum_{a \in (\ZZ/q\ZZ)^\times}
  \Bigl(V_N(a) - \frac{N}{q{-}1}\Bigr)^2.
\]
\end{theorem}

This establishes the equivalence
\defname{SVE}~$\Leftrightarrow$~\defname{VE}, via the identity
$\mathrm{excessEnergy} = (p{-}1)\sum_a (V_a - N/(p{-}1))^2$
(\lean{EM/Reduction/Master.lean}{sve\_controls\_visit\_deviation}): subquadratic visit
energy (the total non-trivial character energy is $o(N^2)$) is the same
as visit equidistribution (the visit counts are approximately uniform).
The identity connects the Fourier-analytic perspective (character sums)
with the geometric perspective (visit counts) on the walk:
\begin{itemize}[nosep]
\item \textbf{SVE} says the walk has no sustained spectral bias:
  projecting onto any non-trivial character gives cancellation.
\item \textbf{VE} says the walk has no sustained spatial bias:
  it visits every position approximately equally often.
\end{itemize}
These are dual formulations of the same condition, connected by
Parseval's identity on the finite group $(\ZZ/q\ZZ)^\times$.

The fiber decomposition (\S\ref{sec:fiber-energy}) refines this further:

\begin{theorem}[\lean{EM/Reduction/VisitEquidist.lean}{fiberMultCharSum\_step}]
The fiber character sum $F(a, N) = \sum_{\substack{n < N \\ w(n) = a}}
\chi(m(n))$ satisfies the recurrence:
$F(a, N{+}1) = F(a, N) + \chi(m(N)) \cdot \mathbf{1}_{w(N) = a}$.
\end{theorem}

\noindent This step-by-step evolution of fiber sums provides the
infrastructure for the CME~$\Rightarrow$~CCSB proof via energy
tracking: the total character energy decomposes into active-fiber
contributions at each step, and CME bounds each fiber's contribution.

\begin{remark}[Almost-all MC and the multiplier/walk gap]
\label{rem:almost-all-mc}
A natural question: since SD gives character sum cancellation for
\emph{almost all} squarefree starting points
(Theorem~\ref{thm:sd-cancel}), does this yield
$\mathrm{GenMullinConjecture}(n)$ for almost all~$n$?

It does not.  The cancellation that SD produces is for the
\textbf{multiplier} character energy
$\|\sum_{k<K} \chi(\seq{n}(k) \bmod q)\|^2 = o(K^2)$---the
individual multiplier residues have cancelling character sums.  But
the Weyl Hitting Bridge (Theorem~\ref{thm:weyl-bridge}), which
converts character cancellation into cofinal hitting, requires
\textbf{walk} character cancellation:
$\|\sum_{k<K} \chi(\walkZ{n}{q}{k})\|^2 = o(K^2)$.  The walk is a
partial product of the multipliers---$\walkZ{n}{q}{k} = \prod_{j<k}
\multZ{n}{q}{j}$---and cancellation of a sum $\sum \chi(m_k)$ does not
imply cancellation of $\sum \prod_{j<k} \chi(m_j)$.

This is Dead End~\#117 (MultCancelToWalkCancel, \S\ref{sec:oracle}):
an explicit counterexample (alternating multipliers $\{2,3\} \bmod 5$
with $\chi(2) = i$) achieves multiplier cancellation yet
$|W_K| = \Theta(K)$, and is compatible with all EM structural
properties.  The chain therefore breaks as:
\[
  \mathrm{SD}
  \;\xrightarrow{\text{proved}}\;
  \text{a.a.\ multiplier cancellation}
  \;\xrightarrow{\text{\#117}}\;
  \text{a.a.\ walk cancellation}
  \;\xrightarrow{\text{proved}}\;
  \text{a.a.\ GenMC}.
\]
The last arrow composes the Weyl bridge with
Theorem~\ref{thm:gen-mc}; the first arrow is
Theorem~\ref{thm:sd-cancel}; the middle arrow is the
marginal/joint barrier, which persists orbit by orbit.
Ensemble averaging helps with multiplier statistics but cannot cross
this barrier: the walk character sum is a \emph{joint} statistic
(depending on all previous multipliers through the partial product),
while the multiplier character sum is a \emph{marginal} one.

This is why the orbit hypothesis is CME rather than multiplier cancellation:
CME goes directly from conditional multiplier equidistribution to
walk cancellation (via the fiber decomposition and telescoping),
bypassing MultCancelToWalkCancel entirely.

Everything above is a statement about \emph{this} route.  A different
route does reach almost-all capture, unconditionally: the seed-average law
of \S\ref{sec:seed-average}, which for each fixed prime~$q$ and each
$\varepsilon>0$ produces a finite horizon~$n$ at which the seeds coprime
to~$q$ missing~$q$ in their first~$n$ multipliers have upper natural
density at most~$\varepsilon$
(Theorem~\ref{thm:almost-all-capture},
\lean{EM/Population/AlmostAllDensity.lean}{almost\_all\_genmc\_density}).
It gets there precisely by \textbf{avoiding character sums entirely}:
capture is converted into membership of one residue in an explicitly
enumerated set (\lean{EM/Population/SeedCapture.lean}{captured\_iff\_mem\_visited}),
after which the argument is exact CRT counting against a deterministic
compensator bound.  The multiplier/walk gap of this remark ---
Dead End~\#117 --- is therefore not on its path, since no walk character
sum is ever formed.  What that route does \emph{not} do is say anything
about the orbit of~$2$: a density-one statement over the seed ensemble is
silent about any single trajectory, and Dead End~\#90 is untouched.
\end{remark}

\subsection{Self-Correcting Feedback and the Lyapunov Route}
\label{sec:self-correcting}

The routes of Sections~\ref{sec:ensemble}
and~\ref{sec:character} target CME (conditional multiplier
equidistribution) or its equivalents.  A structurally different route
bypasses CME entirely: prove \textbf{visit equidistribution}~(VE)
directly, and let the spectral identity of \S\ref{sec:visit-eqd} do
the rest.

Define the empirical frequency vector
\[
  f_c(N) \;=\; \frac{1}{N}\,\bigl|\{n < N : \walkZ{2}{q}{n} = c\}\bigr|,
  \qquad c \in (\ZZ/q\ZZ)^\times,
\]
so that $f_c(N) = V_N(c)/N$ in the notation of the spectral identity
(Theorem~\ref{thm:excess-energy}).
VE asserts $f_c(N) \to 1/\varphi(q)$ for every~$c$.  By the
Weyl criterion this is equivalent to walk character cancellation
$\|\sum_{n<N} \chi(\walkZ{2}{q}{n})\|^2 = o(N^2)$, and the chain
\[
  \text{VE}
  \;\xleftrightarrow[\text{proved}]{\text{Weyl}}\;
  \text{WalkCharCancel}
  \;\xrightarrow[\text{proved}]{\text{WeylHittingBridge}}\;
  \text{MC}
\]
is fully machine-verified
(\code{sve\_controls\_visit\_deviation}, \code{cancel\_weyl\_implies\_mc}).

\medskip
\noindent\textbf{VE is strictly different from CME.}\enspace
CME controls the \emph{conditional} distribution of multipliers given
walk position; VE controls the \emph{marginal} distribution of walk
positions.  Neither implies the other: equidistributed walk positions
do not force equidistributed conditional multipliers, and vice versa.
This is not an equivalence collapse.

\medskip
\noindent\textbf{The self-correcting feedback mechanism.}\enspace
The EM walk has a built-in negative feedback loop.  Suppose $f_c(N)$
is anomalously large: the walk visits position~$c$ too often.
At each such visit, $\Prod{2}(n) \equiv c \pmod{q}$, so
$\Prod{2}(n) + 1 \equiv c + 1 \pmod{q}$.  The next multiplier is
$\seq{}{n{+}1} = \minFac(\Prod{2}(n) + 1)$, an integer whose residue
class mod~$q$ is determined by $\Prod{2}(n) + 1$.  Since $\Prod{2}(n)+1$
is coprime to the bag, its least prime factor is a prime outside the
bag, and once the bag has swallowed the small primes those factors are
large and their residues mod~$q$ have no reason to prefer a class.
(An earlier draft invoked ``Alladi's theorem'' here as unweighted
equidistribution of the least prime factor of generic integers; that
statement is false---Dead End~\#160---and the heuristic is the
orbit-level one just given.)  If the multiplier at position~$c$ is
approximately equidistributed, the transition $c \mapsto c \cdot m(n)$ scatters the
walk away from~$c$ toward all other positions---a negative feedback
that opposes further concentration at~$c$.

This has been formalized via a Lyapunov function
(\code{EM/Reduction/SelfCorrecting.lean}).  Define
\[
  L(N) \;=\; \sum_{c \in (\ZZ/q\ZZ)^\times}
  \Bigl(V_c(N) - \frac{N}{\varphi(q)}\Bigr)^{\!2},
\]
where $V_c(N) = |\{n < N : \walkZ{2}{q}{n} = c\}|$ is the visit count.
The one-step algebraic identity
(\lean{EM/Reduction/SelfCorrecting.lean}{lyapunov\_one\_step})
\[
  L(N{+}1) = L(N) + 2\,d_{\walkZ{2}{q}{N}}(N) + \frac{q-2}{q-1},
\]
where $d_c(N) = V_c(N) - N/(q{-}1)$ is the visit deviation at
position~$c$, telescopes to
(\lean{EM/Reduction/SelfCorrecting.lean}{lyapunov\_telescope})
\[
  L(N) = 2\,R(N) + N\,\frac{q-2}{q-1},
  \qquad R(N) = \sum_{n<N} d_{\walkZ{2}{q}{n}}(n).
\]
The cumulative drift $R(N)$ records how often the walk visits
positions that are already over-represented.  Under a uniform
multiplier measure, $\mathbb{E}[d_c] = 0$ for every position
(\lean{EM/Reduction/SelfCorrecting.lean}{uniform\_multiplier\_zero\_drift}),
and the transition matrix is doubly stochastic
(\lean{EM/Reduction/SelfCorrecting.lean}{group\_walk\_doubly\_stochastic}),
so the drift has zero bias.
\textbf{SelfCorrectingDrift}~(SCD) asserts $R(N) = o(N^2)$, which
gives $L(N) = o(N^2)$, hence VE, and therefore MC via the proved
chain (\lean{EM/Reduction/SelfCorrecting.lean}{scd\_implies\_mc}).

\medskip
\noindent\textbf{SCD $\equiv$ SVE (Dead End~\#120).}\enspace
The Lyapunov telescope $L(N) = 2R(N) + O(N)$ shows that SCD
($R(N) = o(N^2)$) is algebraically equivalent to SVE ($L(N) = o(N^2)$)
for each fixed prime~$q$ above the threshold.  The converse is also
proved: SVE implies SCD above the SVE threshold
(\lean{EM/Reduction/SelfCorrecting.lean}{sve\_implies\_scd\_above\_threshold}).
Therefore SCD provides no new proof leverage beyond SVE; it is a
Lyapunov reformulation that clarifies the dynamical mechanism
(self-correcting feedback) but does not reduce the mathematical
difficulty.

\medskip
\noindent\textbf{What remains.}\enspace
The residual difficulty is orbit-specificity: $\Prod{2}(n) + 1$ is a
specific value, not a random draw from its residue class, and no
population statement supplies its least prime factor's residue.  The Lyapunov analysis shows that proving
$R(N) = o(N^2)$ requires controlling how often the walk visits
positions that are already over-represented---and this control depends
on the specific EM orbit, not just generic multiplier statistics.
This is the same orbit-specificity barrier
from~\S\ref{sec:landscape}, now expressed at the level of single-step
drift rather than long-range character cancellation.

\medskip
\noindent\textbf{Per-class bad-set escape (Dead End~\#121).}\enspace
A natural attempt: use SMSB (SecondMomentSquaredBound) to bound the
\emph{global} density of ``bad'' steps (those whose tail has large
character energy), then argue that each walk-position class has
proportionally few bad steps, then use SubgroupEscape to get escape
from each class.  The per-class bad sets partition the total bad set
(\lean{EM/Reduction/SMSB.lean}{badVisitsInClass\_card\_sum\_eq}),
and pigeonhole guarantees \emph{some} class with small bad density
(\lean{EM/Reduction/SMSB.lean}{exists\_class\_small\_bad\_set}).
But MC requires the specific class~$-1$, not an existential over
classes.  Targeting~$-1$ requires that bad steps do not concentrate
in one class---a per-fiber density statement
(BadSetEquidistribution) that is a consequence of CME\@.  Going from
the marginal bad-density bound (what SMSB provides) to the conditional
bound (what per-class escape requires) is blocked by the
Marginal/Joint Barrier: the joint distribution of (walk position,
badness) is controlled by CME, not by SMSB alone.

\medskip
\noindent\textbf{The LinearDrift weakening.}\enspace
A weaker variant of SCD also suffices: \emph{LinearDrift}
($|R(N)| \leq CN$) implies MC
(\lean{EM/Reduction/WeakRecurrence.lean}{linearDrift\_implies\_mc})
via a threshold-free route:
LinearDrift~$\Rightarrow$~VisitGrowth~$\Rightarrow$~HH%
~$\Rightarrow$~MC, where VisitGrowth ($V(a,N) \to \infty$ for
all~$a$) bypasses the $\mathrm{FiniteMCBelow}$ threshold entirely.
A free unconditional lower bound
$R(N) \geq -N(q{-}2)/(2(q{-}1))$ is proved
(\code{cumulativeDrift\_lower\_bound}), so only the one-sided upper
bound $R(N) \leq CN$ is needed.

\subsection{Tail Window Decorrelation}
\label{sec:twd}

The tail identity (Theorem~\ref{thm:tail-id}) connects temporal
blocks of the standard EM orbit to the ensemble framework: the
character sum over steps $jK{+}1, \ldots, (j{+}1)K$ of the standard
orbit equals the generalized partial sum starting from $\Prod{2}(jK)$.
This motivates decomposing the full character sum into non-overlapping
\emph{windows}.

\begin{definition}[\defname{WindowCharSum}]
The character sum over the $j$-th window of length~$K$:
$S_j = \sum_{k < K} \chi(\seq{2}(jK{+}k{+}1) \bmod q)$.
\end{definition}

\begin{theorem}[\lean{EM/Reduction/TailWindow.lean}{windowCharSum\_eq\_genSeqCharPartialSum}]
$S_j$ equals the generalized character partial sum
starting from $\Prod{2}(jK)$:
$S_j = \sum_{k<K} \chi(\text{genSeq}(\Prod{2}(jK), k) \bmod q)$.
\end{theorem}

\begin{theorem}[\lean{EM/Reduction/TailWindow.lean}{charSum\_eq\_window\_sum}]
\label{thm:window-decomp}
The total character sum decomposes:
$\sum_{k < JK} \chi(\seq{2}(k{+}1) \bmod q) = \sum_{j<J} S_j$.
\end{theorem}

Applying the variance identity
(\lean{EM/Reduction/TailWindow.lean}{norm\_sq\_sum\_eq\_diag\_plus\_cross}),
the total character energy decomposes as
\[
  \Bigl\|\sum_{j<J} S_j\Bigr\|^2
  = \sum_{j<J} |S_j|^2
  + 2 \sum_{j < l < J} \mathrm{Re}(S_j \cdot \overline{S_l}).
\]
The diagonal terms are bounded by $J \cdot K^2$
(\lean{EM/Reduction/TailWindow.lean}{diag\_energy\_le}).
The \emph{cross terms}---the off-diagonal correlations between
windows at different temporal positions---are the new quantity of
interest.

\begin{property}[\defname{TailWindowDecorrelation} (\abbr{TWD})]
The average absolute cross term between non-overlapping windows
vanishes as the window size~$K$ grows:
\[
  \frac{1}{J^2} \sum_{j < l < J}
  |\mathrm{Re}(S_j \cdot \overline{S_l})| \to 0.
\]
\end{property}

\noindent\emph{As stated, TWD is false} (Dead End~\#156,
\lean{EM/Ensemble/UncenteredRefutations.lean}{not\_tailWindowDecorrelation}):
it quantifies over all $\chi$ with $\|\chi\| \le 1$, and for
$\chi \equiv 1$ two windows of length~$K$ have cross term~$K^2$; the
bridge \code{TWDImpliesCCSB} therefore holds vacuously
(\code{twdImpliesCCSB\_vacuous}) and is retired as an open point.  The
centered version---nontrivial Dirichlet characters---is the intended
temporal statement, and the discussion below is about it.

\begin{theorem}[\lean{EM/Reduction/TailWindow.lean}{twd\_implies\_mc\_from\_bridge}]
$\mathrm{TWDImpliesCCSB} + \mathrm{TWD}
\;\Longrightarrow\; \MC$.
\end{theorem}

\noindent The proof composes the open bridge TWDImpliesCCSB with the
proved chain CCSB~$\Rightarrow$~MC\@.  The bridge TWDImpliesCCSB
would follow from TWD by choosing $K = \lceil 1/\varepsilon\rceil$:
the diagonal contributes $N \cdot K = O(\varepsilon^{-1} N)$ while the
cross terms contribute $O(\varepsilon N^2)$, giving
$\|S(N)\|^2 = O(\varepsilon N^2)$ for all~$\varepsilon > 0$, hence CCSB\@.

\begin{remark}[TWD vs.\ ensemble decorrelation]
\label{rem:twd-vs-ensemble}
TWD is a \emph{temporal} decorrelation statement about the standard EM
orbit: windows far apart in the sequence should have nearly
independent character sums.  The physical intuition is that the
accumulator $\Prod{2}(jK)$ grows super-exponentially, so the starting
points of consecutive windows are extremely different in magnitude,
making it hard for any fixed modulus~$q$ to see systematic correlation
between them.

This contrasts with the \emph{fiber} decorrelation used in the
variance route (\S\ref{sec:variance-route}): there, one averages over
squarefree starting points at two fixed time steps $j, k$.  TWD
averages over time windows at the single starting point $n = 2$.  The
two approaches face different barriers: the variance route faces the
orbit-specificity barrier (density-one results cannot pin down $n=2$),
while TWD faces the temporal-correlation barrier (consecutive windows
share the accumulator prefix $\Prod{2}(jK)$).

The file \code{EM/Reduction/TailWindow.lean} (477~lines, zero
\code{sorry}) formalizes the window definitions, the tail identity
connection, the variance identity, the diagonal and cross-term bounds,
and the conditional chain TWD~$\Rightarrow$~CCSB~$\Rightarrow$~MC\@.
\end{remark}
